\input{estilo}
% El cómputo A con dos nodos, cada uno con su rango y su watermark; los dos le
% mandan datos a B y le informan su watermark a la autoridad, que calcula el
% mínimo y se lo pasa a B. Los watermarks en rojo, como en el resto de la clase.
\begin{document}
\begin{figura}\begin{tikzpicture}
  \node[caja] (a1) at (0,1.2) {A};
  \node[caja] (a2) at (0,-1.2) {A};
  \node[etiq, left=8pt of a1] {[a–m]};
  \node[etiq, left=8pt of a2] {[m–z]};
  \node[etiq, text=acento, above=2pt of a1] {\var{W} = 100};
  \node[etiq, text=acento, above=2pt of a2] {\var{W} = 200};
  \draw (0,0) ellipse[x radius=1.75, y radius=2.35];
  \node[caja] (b) at (4.2,0.6) {B};
  \node[caja, minimum width=28mm, minimum height=10mm, align=center] (au) at (4.2,-2.6) {autoridad de\\watermarks};
  \draw[msg] (a1.east) -- (b.west);
  \draw[msg] (a2.east) -- (b.west);
  \draw[msg, draw=suave] (a1.south east) to[out=-30, in=135] (au.north west);
  \draw[msg, draw=suave] (a2.south east) to[out=-50, in=180] (au.west);
  \draw[msg, draw=acento] (au.north) -- node[etiq, text=acento, right] {\var{W}(A) = 100} (b.south);
  \node[etiq, align=left, anchor=west] (t) at ($(au.east)+(0.3,0)$) {A: [a–m] = 100\\\phantom{A: }[m–z] = 200};
  \draw[decorate, decoration={brace, amplitude=5pt}] (t.north east) -- (t.south east);
  \node[etiq, anchor=west, text=acento] at ($(t.east)+(0.3,0)$) {\ensuremath{\Rightarrow} A = 100};
\end{tikzpicture}\end{figura}
\end{document}
